% Part: normal-modal-logic
% Chapter: sequent-calculus
% Section: introduction

\documentclass[../../../include/open-logic-section]{subfiles}

\begin{document}

\olfileid{nml}{seq}{int}

\olsection{પરિચય}

વિધાનતર્કના સિક્વન્ટ કલનમાં
\iftag{prvBox}{$\Box$\iftag{prvDiamond}{ અને~}{}}{}%
\iftag{prvDiamond}{$\Diamond$}{} માટેના વધારાના નિયમો
ઉમેરી શકાય છે. દાખલા તરીકે,~$\Log{K}$ માટે
\Log{LK} ઉપરાંત આ નિયમો છે:
\[
  \iftag{prvBox}{\iftag{prvDiamond}{% <> and [] primitive
    \Axiom$\Gamma \fCenter \Delta, !A$
    \RightLabel{$\Box$}
    \UnaryInf$\Box\Gamma \fCenter \Diamond\Delta, \Box!A$
    \DisplayProof  
    \qquad
    \Axiom$!A, \Gamma \fCenter \Delta$
    \RightLabel{$\Diamond$}
    \UnaryInf$\Diamond!A, \Box\Gamma \fCenter \Diamond\Delta$
    \DisplayProof
}{% only Box primitive
\Axiom$\Gamma \fCenter !A$
\RightLabel{$\Box$}
\UnaryInf$\Box\Gamma \fCenter \Box!A$
\DisplayProof
}}{% only <> primitive
  \Axiom$!A \fCenter \Delta$
  \RightLabel{$\Diamond$}
  \UnaryInf$\Diamond!A \fCenter \Diamond\Delta$
  \DisplayProof
}
\]
$\Log{K}$ના વિસ્તરણો માટે બીજા નિયમો પણ ઉમેરવા પડે છે.

દરેક મોડલ તર્ક માટે આવું સિક્વન્ટ કલન જાણીતું નથી.
$\Log{S5}$ અર્થવિજ્ઞાનની દૃષ્ટિએ સરળ છે (તેને સુલભતા
સંબંધ વાપર્યા વિના પણ વ્યાખ્યાયિત કરી શકાય છે), છતાં
$\Log{LK}$માંથી મળતું અને~\Cut નિયમ વિના પૂર્ણ હોય એવું
સિક્વન્ટ કલન તેના માટે જાણીતું નથી. જોકે તેના માટે કટ વિના
પૂર્ણ \emph{હાઇપર-સિક્વન્ટ} કલન છે.

\end{document}
